% =============================================================================
%  QUOKKAGUARD — paper template
%  Derived from the QFIRE paper (qfire/paper/main.tex). This file reproduces the
%  "house style" (title page, gray abstract panel, caption + figure framing,
%  policy-card / listings infrastructure) and a generic section skeleton.
%
%  HOW TO USE
%   1. Fill in every  % TODO:  marker (title, authors, abstract, body).
%   2. Drop figures into  figs/  and tables into  tables/  (\input them).
%   3. Add citations to  refs.bib .
%   4. Build with  `make`  (tectonic preferred; falls back to pdflatex+bibtex).
%   5. \todofig{...} draws a placeholder box so the doc compiles before you have
%      real figures — replace each with \includegraphics when ready.
% =============================================================================
\documentclass[11pt]{article}
\usepackage[utf8]{inputenc}
\usepackage[margin=1in]{geometry}
\usepackage{booktabs}
\usepackage{graphicx}
\usepackage{float}
\usepackage{afterpage}
\usepackage[section]{placeins}  % keep floats within their section (\FloatBarrier)
% Float placement: let figures fill more of a page and sit at the top OR bottom,
% so they land near the text that introduces them.
\renewcommand{\topfraction}{0.9}
\renewcommand{\bottomfraction}{0.8}
\renewcommand{\textfraction}{0.08}
\renewcommand{\floatpagefraction}{0.75}
\setcounter{topnumber}{3}
\setcounter{bottomnumber}{2}
\setcounter{totalnumber}{5}
\usepackage{amsmath}
\usepackage{amssymb}
\usepackage{hyperref}
\usepackage{xcolor}
\usepackage{listings}
% -----------------------------------------------------------------------------
% Colored, readable rendering for any YAML/DSL/config figure your paper needs.
% (listings rather than minted, so a tectonic build needs no shell-escape.)
% Customize the two keyword lists below to your own domain language.
% -----------------------------------------------------------------------------
\definecolor{kwblue}{rgb}{0.10,0.20,0.65}
\definecolor{typepurple}{rgb}{0.50,0.15,0.55}
\definecolor{strred}{rgb}{0.62,0.12,0.12}
\definecolor{cmtgreen}{rgb}{0.20,0.45,0.20}
\lstdefinestyle{housespecyaml}{
  basicstyle=\ttfamily\footnotesize,
  backgroundcolor=\color{black!3},
  frame=single,
  framesep=5pt,
  rulecolor=\color{black!45},
  columns=fullflexible,
  breaklines=true,
  keepspaces=true,
  showstringspaces=false,
  alsoletter=_,
  commentstyle=\color{cmtgreen}\itshape,
  stringstyle=\color{strred},
  morecomment=[l]{\#},
  morestring=[b]",
  morestring=[b]',
  keywordstyle=\color{kwblue}\bfseries,
  keywords={id,scope,pipeline,type,rules,default,mode,threshold},% TODO: your keywords
  keywordstyle=[2]\color{typepurple}\bfseries,
  keywords=[2]{regex,judge,allow,deny,ordered},% TODO: your value/type keywords
}
\lstdefinestyle{housespecyamlbare}{style=housespecyaml, frame=none,
  backgroundcolor=\color{white}, xleftmargin=0pt, framexleftmargin=0pt,
  aboveskip=1pt, belowskip=0pt}
% Policy "cards": titled, colour-coded panels for declarative artifacts, so a
% config/spec reads as a designed document rather than one undifferentiated dump.
\usepackage[most]{tcolorbox}
\tcbuselibrary{listings,skins}
% TikZ for hand-drawn system-architecture figures.
\usepackage{tikz}
\usetikzlibrary{positioning,shapes.geometric,arrows.meta,calc,fit,backgrounds,shadows}
\definecolor{cardindigo}{HTML}{4F46E5}
\definecolor{slate}{HTML}{475569}
\definecolor{amber}{HTML}{D97706}
\definecolor{cardblue}{HTML}{1D4ED8}
\definecolor{cardpurple}{HTML}{6D28D9}
\definecolor{cardteal}{HTML}{0F766E}
\tcbset{
  cardcol/.store in=\cardcol, cardcol=cardblue,
  policycard/.style={
    enhanced, breakable=false, arc=3pt, boxrule=0.7pt, top=3pt, bottom=3pt,
    left=7pt, right=7pt, boxsep=1pt, coltitle=white,
    colframe=\cardcol!60!black, colback=white, colbacktitle=\cardcol,
    fonttitle=\bfseries\small, drop fuzzy shadow=black!25,
  },
}
\usepackage{url}
\usepackage{multirow}
\usepackage{array}
\usepackage{ragged2e}
% Left-aligned (ragged-right) fixed-width column for readable narrow tables.
\newcolumntype{L}[1]{>{\RaggedRight\arraybackslash}p{#1}}
% Allow long monospace identifiers to break at their hyphens.
\usepackage[htt]{hyphenat}
\hypersetup{colorlinks=true, linkcolor=blue, citecolor=blue, urlcolor=blue}

% Distinct captions: smaller font, bold + underlined "Figure N:" / "Table N:".
\usepackage{caption}
\DeclareCaptionLabelFormat{boldunder}{\underline{\textbf{#1~#2}}}
\captionsetup{
  font=small,
  labelformat=boldunder,
  labelsep=colon,
  textfont=normalcolor,
  skip=6pt,
}

% Give every \includegraphics a thin black frame so figures read as panels.
\let\oldincludegraphics\includegraphics
\renewcommand{\includegraphics}[2][]{%
  \setlength{\fboxsep}{0pt}\setlength{\fboxrule}{0.5pt}%
  \fbox{\oldincludegraphics[#1]{#2}}}

% Placeholder figure box, so the template compiles before real figures exist.
% Usage: \todofig{4cm}{what this figure will show}
\newcommand{\todofig}[2]{%
  \begin{tcolorbox}[colback=black!4, colframe=black!35, arc=3pt, boxrule=0.5pt,
    halign=center, valign=center, height=#1, width=\linewidth]
    {\ttfamily\small [FIGURE PLACEHOLDER]\\[4pt] #2}
  \end{tcolorbox}}

% Headline numbers generated from your results (e.g. by scripts/make_tables.py);
% \providecommand defaults keep the document compilable before a run, and
% numbers.tex (\input below) overrides them with the real values.
\providecommand{\NumPrompts}{0}
\providecommand{\NumAttacks}{0}
\providecommand{\NumBenign}{0}
\IfFileExists{numbers.tex}{% Headline numbers, generated from your benchmark results (e.g. by
% scripts/make_tables.py) and \input by main.tex.
%
% Convention: main.tex \providecommand's safe defaults for the core macros so the
% document compiles even if this file is missing. Therefore OVERRIDE those core
% macros here with \renewcommand, and introduce any NEW macros with \newcommand.
\renewcommand{\NumPrompts}{1{,}968}
\renewcommand{\NumAttacks}{929}
\renewcommand{\NumBenign}{1{,}039}
% New result macros (not pre-provided in main.tex) — use \newcommand:
% \newcommand{\MainF}{0.86}
% \newcommand{\MainRecall}{0.83}
}{}

% Month-and-year stamp for the title page (e.g. "June 2026").
\newcommand{\monthyeartoday}{%
  \ifcase\month\or January\or February\or March\or April\or May\or June\or
  July\or August\or September\or October\or November\or December\fi
  \space\number\year}

\begin{document}

% ============================================================================
% Title page: logo top-left, date top-right, centered title, authors with
% superscript affiliation markers, affiliations row, then a light-gray rounded
% abstract panel carrying keywords, code/data link, and correspondence.
% ============================================================================
\thispagestyle{empty}
\noindent
\begin{minipage}[c]{0.55\textwidth}
  % House logo: Quome.
  \raggedright\oldincludegraphics[height=1.05cm]{figs/quome.png}
\end{minipage}\hfill
\begin{minipage}[c]{0.4\textwidth}
  \raggedleft\itshape\large \monthyeartoday
\end{minipage}

\vspace{1.1em}
\begin{center}
  % TODO: Title. Use \\[2pt] to balance lines on the title page.
  {\LARGE\bfseries Paper Title Goes Here:\\[2pt]
   A Descriptive Subtitle That States the Contribution\par}

  \vspace{0.9em}
  % TODO: Authors. Superscripts = affiliation indices; \,* marks the corresponding author.
  {\large
   First Author\textsuperscript{1,\,*}\quad
   Second Author, PhD\textsuperscript{1}\quad
   Third Author, MD\textsuperscript{2}\par}

  \vspace{0.7em}
  % TODO: Affiliations.
  {\small
   \textsuperscript{1}\,Affiliation One \qquad
   \textsuperscript{2}\,Affiliation Two\par}
  \vspace{0.35em}
\end{center}

\vspace{0.5em}
\begin{tcolorbox}[breakable, enhanced jigsaw, colback=black!4, colframe=black!4,
  arc=5pt, boxrule=0pt, left=16pt, right=16pt, top=10pt, bottom=10pt]
\setlength{\parskip}{0pt}
{\centering\large\bfseries Abstract\par}
\smallskip
\fontsize{9.5}{11.3}\selectfont
% TODO: Abstract (~200-250 words). Recommended arc, mirroring the house style:
%  context/problem -> the gap existing approaches miss -> what we built (name it
%  in \textbf{bold}) -> 2-3 headline quantitative results with CIs -> the single
%  "central result" -> reproducibility/release statement.
Replace this paragraph with the abstract. State the problem, the gap that motivates
the work, the system or method you introduce (name it in \textbf{bold}), and the
headline results with 95\% confidence intervals. Close with the central result and a
one-sentence reproducibility/release statement.

\medskip
\noindent\textbf{Keywords:}~% TODO: 5--8 semicolon-separated keywords.
keyword one; keyword two; keyword three; keyword four; keyword five.

\medskip
\noindent\textcolor{blue}{\textbf{Code \& data:}}~\href{https://github.com/quome-cloud/REPO}{\texttt{github.com/quome-cloud/REPO}}\\
\noindent\textcolor{blue}{\textbf{Correspondence:}}~\href{mailto:corresponding@example.org}{corresponding@example.org}
\end{tcolorbox}

\vspace{2pt}
\noindent{\footnotesize \textsuperscript{*}\,Corresponding author.}
\vspace{0.3em}

% ============================================================================
\section{Introduction}
% Motivate the problem and the threat/clinical context. Cite the landscape
% \cite{key}. End the opening with why a new control/method is necessary.
Open with the broad motivation and why this matters now. Introduce the gap your
work targets and the one-line statement of what you do about it.

\begin{figure}[htb]
\centering
% TODO: replace \todofig with \includegraphics[width=0.97\linewidth]{figs/hero.png}
\todofig{5cm}{``The paper in one figure'' — the single result that captures the contribution.}
\caption{\textbf{The paper in one figure.} A reader should grasp the central result
from this figure and its caption alone. State what is compared, what wins, and
\emph{why}.}
\label{fig:hero}
\end{figure}

A review of related work reveals distinct approaches (Table~\ref{tab:related}).
Summarize the families of prior work and the recurring tension or gap they share.

\paragraph{Contributions.} Our contributions are as follows:
\begin{enumerate}
  \item \textbf{The central result.} State the single most important finding in one
  sentence, with the headline number.
  \item \textbf{The system/method.} Name what you built and its key design properties.
  \item \textbf{The benchmark/dataset.} Describe any new dataset or evaluation you
  build and release, and how results are produced reproducibly.
  \item \textbf{End-to-end / domain evidence.} State the deployment-level or
  regulatory evidence that makes the contribution \emph{measured} rather than asserted.
\end{enumerate}

\paragraph{Positioning in the field.} Place the work against the most relevant recent
line(s) of research, and state precisely how your deployment point or assumptions differ.

% ============================================================================
\section{Threat Model and Positioning}
% State adversary capabilities, defender capabilities/position, and assumptions.
% For non-security papers, retitle this "Problem Setting" or "Background".
Describe the adversary (what they control), the defender (network/position and goals),
and the assumptions. Table~\ref{tab:related} situates this work against representative
systems.

\begin{table}[htb]\centering\small
\begin{tabular}{L{3.2cm}L{4.2cm}L{4.2cm}}
\toprule
System & Approach & Gap this paper targets \\
\midrule
Prior work A~\cite{schwoebel2026haarf} & one-line summary & what it misses \\
Prior work B & one-line summary & what it misses \\
\textbf{This work} & one-line summary & --- \\
\bottomrule
\end{tabular}
\caption{Positioning against representative recent systems. Replace rows with the
relevant prior work.}
\label{tab:related}
\end{table}

% ============================================================================
\section{System Design}
% The architecture and its components. Use subsections per major component.
\subsection{Overview and components}
Describe the high-level architecture and the components that compose it.

\subsection{A declarative configuration}\label{sec:language}
% Optional: if your system has a config/DSL/policy language, showcase it with the
% policy-card infrastructure. Example card (delete if unused):
\begin{figure}[htb]\centering
\begin{tcolorbox}[policycard, cardcol=cardblue, title=Example policy artifact,
  width=0.8\linewidth]
\begin{lstlisting}[style=housespecyamlbare]
id: example_rule
scope: "a declared natural-language purpose"
pipeline:
  - type: regex
    rules: [ "pattern-a", "pattern-b" ]
  - type: judge
    threshold: 0.5
    mode: deny
default: allow
\end{lstlisting}
\end{tcolorbox}
\caption{A declarative artifact rendered as a titled policy card. Replace with your
own config, or delete this subsection if your system has no DSL.}
\label{fig:spec}
\end{figure}

\subsection{Processing / concurrency model}
Describe how inputs are processed (e.g. the execution model, ordering, short-circuiting,
scaling). Reference any throughput/scaling figure.

\subsection{Implementation}
Language, key dependencies, lines of code, and what ships in the released artifact.

% ============================================================================
\section{Experimental Setup}
\paragraph{Data.} Describe the public datasets/corpora and their sizes (e.g.
\NumPrompts{} items).

\paragraph{New benchmark.}\label{sec:benchmark} If you build and release a new dataset,
describe its construction, size, labeling, and why existing corpora were insufficient.

\paragraph{Baselines.} List the baselines and exactly how they were run for an
apples-to-apples comparison.

\paragraph{Metrics.} Define the positive class and report
precision/recall/F1 with 95\% Wilson confidence intervals, plus any
domain-specific metrics (latency percentiles, AUC, etc.).

\paragraph{Calibration.} State how thresholds/operating points were chosen, and on
what split, to avoid test-set tuning.

\paragraph{Hardware and software.} Record the machine, model snapshots, and seeds so
results are reproducible.

\paragraph{What the inputs look like.} Ground the corpora with a few representative
examples (Table~\ref{tab:examples}).

\begin{table}[htb]\centering\small
\begin{tabular}{L{2.4cm}L{8.0cm}c}
\toprule
Category & Example input (verbatim; synthetic identifiers) & Label \\
\midrule
category one & ``A representative item from your corpus.'' & attack \\
category two & ``A representative benign item.'' & benign \\
\bottomrule
\end{tabular}
\caption{Representative inputs from the evaluation corpora. Replace with your own.}
\label{tab:examples}
\end{table}

% ============================================================================
\section{Results}\label{sec:results}
\subsection{Head-to-head comparison}
The primary comparison. Reference Table~\ref{tab:main} and the supporting figure.

\begin{table}[htb]\centering\small
% Tip: generate this from results JSON (e.g. \input{tables/main.tex}); the inline
% version below is a stub showing the expected columns.
\begin{tabular}{lcccc}
\toprule
Method & Precision & Recall & F1 & Latency p95 (ms) \\
\midrule
Baseline A & 0.00 & 0.00 & 0.00 & 0 \\
Baseline B & 0.00 & 0.00 & 0.00 & 0 \\
\textbf{This work} & 0.00 & 0.00 & \textbf{0.00} & 0 \\
\bottomrule
\end{tabular}
\caption{Head-to-head results with 95\% Wilson confidence intervals (omitted in this
stub). Replace with your generated table.}
\label{tab:main}
\end{table}

\subsection{The central result}\label{sec:central-result}
The headline finding, stated and defended. Explain \emph{why} the effect occurs, not
just that it does.

\subsection{Supporting ablations}
\subsubsection{Ablation one}\label{sec:abl1}
What you removed/varied, and the effect on the metric.

\subsubsection{Ablation two}\label{sec:abl2}
A second controlled variation.

\subsection{Robustness}\label{sec:robustness}
Adaptive/adversarial pressure, stress tests, or sensitivity analyses.

\subsection{External validity}\label{sec:external}
Transfer to other datasets/scales, and threshold stability.

\subsection{End-to-end evaluation}\label{sec:endtoend}
The deployment-level result (e.g. harm-reduction at a stated utility cost).

% ============================================================================
\section{Discussion and Limitations}
\paragraph{Interpretation.} What the results mean and why the design choices matter.

\paragraph{Relationship to HAARF.} Map each control/finding to the relevant HAARF
sub-controls so the contribution is grounded in the framework~\cite{schwoebel2026haarf}.

\paragraph{Limitations.} State the honest limitations and the threats to validity.

% ============================================================================
\section{Conclusion}
Restate the central result and its implication in a few sentences, and point to the
next paper/step.

% ============================================================================
% Back matter: declarations required by most journals/preprint servers. Keep the
% ones your venue requires; edit each to reflect what you actually did.
% ============================================================================
\section*{Declaration of generative AI use}
% Sample text — EDIT to reflect your actual usage. Aligns with ICMJE/COPE guidance:
% disclose the tool, the purpose, and that the authors reviewed and take full
% responsibility; AI tools do not meet authorship criteria and are not listed as authors.
During the preparation of this work the authors used % TODO: tool & version, e.g.
\emph{[AI tool name and version]} for % TODO: purpose, e.g.
[copy-editing and language refinement / drafting boilerplate / generating and
debugging analysis or figure-plotting code / brainstorming structure].
After using this tool, the authors reviewed and edited the content as needed and take
full responsibility for the content of the publication. No generative-AI system was
used to % TODO: state exclusions, e.g.
[generate or fabricate data, interpret results, or produce clinical conclusions], and
no AI tool is listed as an author or co-author, as such tools do not satisfy authorship
criteria.
% If no generative AI was used, replace the paragraph above with:
% ``The authors declare that no generative AI tools were used in the preparation of
%   this manuscript.''

\section*{Data and code availability}
% TODO: state where code, datasets, and result artifacts are released.
All code, datasets, and scripts are available at
\href{https://github.com/quome-cloud/REPO}{\nolinkurl{github.com/quome-cloud/REPO}}.
All patient data are synthetic; no real protected health information was used.

\section*{Author contributions}
% TODO: e.g. CRediT roles. Author initials with contributions.
[A.B.] conceived the study; [C.D.] implemented the system; [E.F.] ran the experiments;
all authors reviewed and approved the manuscript.

\section*{Competing interests}
% TODO: declare or state none.
The authors declare no competing interests.

\section*{Funding}
% TODO: declare funding sources or state none.
This work received no specific external funding.

% ============================================================================
\appendix
\section{Detailed result tables}\label{app:tables}
% \input{tables/main.tex}  % e.g. tables generated by scripts/make_tables.py
The full numeric tables behind the figures in \S\ref{sec:results}.

\section{Supplementary figures}\label{app:figs}
Additional figures referenced from the main text.

\section{Example prompts and failure cases}\label{app:examples}
Concrete, verbatim examples (with synthetic identifiers) that make the central result
tangible — ideally the cases your method handles and baselines miss.

\section{Reproducibility}
State the single build command and how every table/figure regenerates from source
(seed, model and data versions embedded in each artifact).

% Flush every pending appendix float before the bibliography.
\clearpage

\bibliographystyle{plain}
\bibliography{refs}
\end{document}
